\documentclass[12pt,a4paper]{report}
\usepackage[utf8]{inputenc}
\usepackage{amsmath}
\usepackage{amsfonts}
\usepackage{amssymb}
\usepackage{amsthm}
\usepackage{hyperref}

\usepackage{multicol}
\usepackage{fancyhdr}
\usepackage[inline]{enumitem}
\usepackage{tikz}
\usepackage{tikz-cd}
\usetikzlibrary{calc}
\usetikzlibrary{shapes.geometric}
\usetikzlibrary{positioning}
\usepackage[margin=0.5in]{geometry}
\usepackage{xcolor}

\hypersetup{
    colorlinks=true,
    linkcolor=blue,
    filecolor=magenta,      
    urlcolor=cyan,
    pdftitle={PDE2<F2>},
    pdfpagemode=FullScreen,
    }

%\urlstyle{same}

\newcommand{\CLASSNAME}{Math 6230 -- Partial Differential Equations II}
\newcommand{\STUDENTNAME}{Paul Carmody}
\newcommand{\ASSIGNMENT}{Extra Credit Assignment \#2}
\newcommand{\DUEDATE}{April 19, 2026}
\newcommand{\PROFESSOR}{Professor Perrera}
\newcommand{\SEMESTER}{Spring 2026}
\newcommand{\SCHEDULE}{MW 11:00 -- 12:15}
\newcommand{\ROOM}{332}

\newcommand{\MMN}{M_{m\times n}}
\newcommand{\FF}{\mathcal{F}}

\pagestyle{fancy}
\fancyhf{}
\chead{ \fancyplain{}{\CLASSNAME} }
%\chead{ \fancyplain{}{\STUDENTNAME} }
\rhead{\thepage}

\fancyfoot{} % clear all footer fields
\fancyfoot[LE,RO]{\thepage}
\fancyfoot[LO,CE]{-------\\\ASSIGNMENT}
%\fancyfoot[CO,RE]{To: Dean A. Smith}
\input{../MyMacros}

\newcommand{\RED}[1]{\textcolor{red}{#1}}
\newcommand{\BLUE}[1]{\textcolor{blue}{#1}}

\newcommand{\SUPP}{\operatorname{supp}}
\newcommand{\CLOSURE}{\operatorname{closure of}}
\newcommand{\BD}{\operatorname{bd}}
\newcommand{\HHN}{\mathcal{H}^n}
\newcommand{\COFF}[3]{\Gamma^{#1}_{{#2}{#3}}}
\newcommand{\VK}[1]{\textbf{#1}}
\newcommand{\VKR}{\VK{R}}
\newcommand{\VKZ}{\VK{Z}}
\newcommand{\CSPACE}{{C(\CONJ{U})}}
\newcommand{\CNORM}[1]{|| #1 ||_\CSPACE}
\newcommand{\CSEMINORM}[1]{[ #1 ]_\CSPACE}
\newcommand{\CSUBSPACE}[3]{C^{#1,#2}(#3)}
\newcommand{\HSPACE}[3]{C^{#1,#2}(#3)}
\newcommand{\HOLDERNORM}[4]{|| #4 ||_{\HSPACE{#1}{#2}{#3}}}
\newcommand{\HSEMINORM}[4]{[ #4 ]_{\HSPACE{#1}{#2}{#3}}}
\newcommand{\SSPACE}[3]{{W^{#1,#2}(#3)}}

\begin{document}

\begin{center}
	\Large{\CLASSNAME -- \SEMESTER} \\
	\large{ w/\PROFESSOR}
\end{center}
\begin{center}
	\STUDENTNAME \\
	\ASSIGNMENT -- \DUEDATE\\
\end{center} 

Let $U \subset \R^n$ be a bound open set and consier the boundary value problem
\begin{align*}
		\BINDEF{-\Delta u + c(x)u = f(x) & \text{in } U}{
		u=0 & \text{on } \partial U}
\end{align*} where $c(x) \ge 0$ for a.e., $x \in U$.  Show that this problem has a unique weak solution $u\in H_0^1(U)$ in each of the following cases:
\begin{enumerate}[label=(\roman*)]
	\item $n=2, c \in L^p(U)$ for some $p>1$, and $f \in L^s(U)$ for some $s >1$.
	\item $n\ge 3, c\in L^{n/2}(U)$, and $f \in L^{2n/(n+2)}(U)$.
\end{enumerate}

\HLINE

\begin{description}

\item Let 
\begin{align*}
	a(u,v) &= \int_U Du\cdot Dv\, dx + \int_U c(x)uv\, dx \AND \\
	L(v) &= \int_U f(x)v(x)\,dx
\end{align*}WTS that $a(u,v) = L(v),\,\forall v \in H_0^1(U)$.

\item \begin{align*}
	c(x) >0 \implies a(u,u) &= \int_U |Du|^2 + \int_U c(x)u^2 \ge \int_U |Du|^2 \\
\end{align*}therefore 
\begin{align*}
	\int_U |Du|^2 \le C\|u\|^2_{H_0^1(U)}
\end{align*}

\item WTS $a(u,v)$ is bounded in each case \begin{description}
	\item Part $(i)$: $n = 2, c \in L^p(U), p >1$.
	\begin{align*}
		H^1_0(U) \hookrightarrow L^q(U),\, \forall q < \infty
	\end{align*}Pick $q$ such that
	\begin{align*}
		p>n\implies \frac{1}{p}+\frac{2}{q}=1
	\end{align*}then due to H\"older
	\begin{align*}
		\int_U cuv \le \|c\|_{L^p}\,\|u\|_{L^q}\,\|v\|_{L^q}	
	\end{align*}Then there exists a $C$ such that
	\begin{align*}
		\|U\|_{L^q}\,\|v\|_{L^p} \le C\|u\|_{H_0^1}\,\|v\|_{H_0^1}.
	\end{align*}Hence
	\begin{align*}
        |a(u,v)| \le C\|u\|_{H_0^1}\,\|v\|_{H_0^1}.
	\end{align*}

	\item Part $(ii)$: $n\ge 3,\,c\in L^{n/2}(U)$.
	
	Using $p*=\frac{2n}{n-2}$ then
	\begin{align*}
		\frac{1}{n/2}+\frac{2}{p*} &= \frac{2}{n}+\frac{2(n-2)}{2n} = \frac{2}{n}+\frac{n-2}{n}=1 \\
		H_0^1(U) &\hookrightarrow L^{p*} 
	\end{align*}consequently
	\begin{align*}
		\int_U cuv \le \|c\|_{L^{n/2}}\,\|u\|_{L^{p*}}\,\|v\|_{L^{p*}}
	\end{align*}Therefore bounded.
    \end{description}

\item WTS that $L(v)$ is bounded in each case \begin{description}

	\item Part $(i)$: $f\in L^s,\, s>1$.

	Let $q$ be such that 
	\begin{align*}
		\frac{1}{s}+\frac{1}{q}=1.
	\end{align*}Since $q$ is finite, it follows that (via Sobolev embedding)
	\begin{align*}
		|L(v)| \le \|f\|_{L^s}\,\|v\|_L^q \le C\|v\|_{H_0^1}.
	\end{align*}

	\item Part $(ii)$: $f \in L^{2n/n+2}$
	
	When we choose $s$ and $q$ as
	\begin{align*}
		s = \frac{2n}{n+2}, q = \frac{2n}{n-2} 
	\end{align*}clearly $\frac{1}{s}+\frac{1}{q} = 1$.  Then applying H\"older's Inequality
	\begin{align*}
		|L(v)| &\le \|f\|_{L^s}\,\|v\|_{L^q} \\
		\AND \|v\|_{L^q} &\le C\|v\|_{H_o^1}
	\end{align*}

\end{description}

\item Consequently, there exists a unique $u \in H_0^1(U)$ such that 
	\begin{align*}
		a(u,v) = L(v),\, \forall H_0^1(U).
	\end{align*}


\end{description}


\end{document}
